\chapter{Boundaries, Defects, And Categories In Topological QFT}
\label{ch:tqft-boundaries-defects-categories}

\ConventionsEuclidean

This chapter develops boundary and defect data in topological QFT.  The
central idea is that a topological theory assigns algebraic objects not only
to closed manifolds, but also to strata, interfaces, and corners, and that
composition of bordisms becomes composition of the corresponding objects.

\section{Boundary Condition Datum}

Let \(Z\) be an \(n\)-dimensional topological theory.  A boundary condition
\(B\) is an assignment of a state space or category to a boundary component
together with compatible amplitudes for bordisms with that boundary label.
For an interval \(I\) with endpoints labelled by \(B_0,B_1\), the theory
assigns a space
\[
  \mathcal H(B_0,B_1).
\]
Gluing two intervals gives a composition map
\[
  \mathcal H(B_0,B_1)\otimes \mathcal H(B_1,B_2)
  \longrightarrow
  \mathcal H(B_0,B_2).
\]
Associativity is the equality of amplitudes for the two decompositions of a
three-segment interval.  Thus boundary conditions in a one-dimensional
topological theory form a category.

\section{Defect Fusion}

A codimension-one defect \(D_{12}\) from a theory \(Z_1\) to a theory \(Z_2\)
is placed on a separating hypersurface.  Fusion is defined by bringing two
parallel defects together:
\[
  D_{12}\circ D_{23}=D_{13}.
\]
For a topological defect, correlation functions are unchanged under
deformations that keep the ordering of defects and avoid operator insertions.
The fusion product is associative when the three-defect collision limit is
independent of the relative distances used to take the limit.  This
associativity is an analytic statement in a QFT realization and an axiom in a
pure bordism-functor presentation.

\section{Line Defects In Three-Dimensional TQFT}

In a three-dimensional oriented TQFT, line defects form a monoidal category
\(\mathcal C\).  The tensor product is parallel fusion:
\[
  X\otimes Y=\lim_{\epsilon\to0}X(\epsilon)Y(0).
\]
Braiding is obtained by exchanging two topological lines in \(\R^3\):
\[
  c_{X,Y}:X\otimes Y\longrightarrow Y\otimes X.
\]
The hexagon identities are equalities of isotopies of three braided lines.
If every line has a dual line and the vacuum line has endomorphism ring
\(\C\), then \(\mathcal C\) is rigid and has a distinguished unit object.
Additional finiteness, semisimplicity, and nondegenerate braiding hypotheses
give a modular tensor category; these hypotheses are part of the claimed
three-dimensional theory, not consequences of topology alone.

\section{Boundary Categories For Chern--Simons Theory}

For compact \(G\) Chern--Simons theory at level \(k\), a Wilson line labelled
by a representation \(R\) is
\[
  W_R(\gamma)=\operatorname{tr}_R\,\mathcal P
  \exp\!\left(\int_\gamma A\right).
\]
Quantum gauge invariance restricts the admissible line labels.  For simply
connected compact simple \(G\), the expected finite label set is the set of
integrable highest weights at level \(k\).  Fusion of line defects is then
encoded by the corresponding representation category of the affine algebra or
quantum group at a root of unity, depending on the construction chosen.

The boundary chiral algebra is not an extra decoration.  It is the algebraic
object that describes states on a surface with boundary and line endpoints.
The equality between bulk Chern--Simons amplitudes and boundary conformal
blocks is a gluing statement relating two constructions of the same
surface-state spaces.

\section{BV-BFV Boundary Mechanism}

For a gauge theory on a manifold \(M\) with boundary, the BV variation has the
form
\[
  \delta S_{\mathrm{BV}}
  =
  \iota_{Q_{\mathrm{BV}}}\omega_{\mathrm{BV}}
  +
  \pi_{\partial M}^*\alpha_{\partial}.
\]
The boundary one-form \(\alpha_{\partial}\) has differential
\[
  \omega_{\partial}=\delta\alpha_{\partial}.
\]
The pair \((\mathcal F_{\partial},\omega_{\partial})\) is the BFV phase
space, and the bulk BV vector field induces a boundary cohomological vector
field when the master equation is compatible with the boundary term.  A
boundary condition is a Lagrangian subspace or derived Lagrangian object
\[
  \mathcal L_{\partial}\subset\mathcal F_{\partial}
\]
on which the boundary variation vanishes.  This is the local mechanism behind
boundary categories in gauge-theoretic TQFTs.

\section{Extended Assignment}

An extended \(n\)-dimensional TQFT assigns:
\[
  Z(M^n)\in\C,\qquad
  Z(\Sigma^{n-1})\in\mathrm{Vect},
\]
\[
  Z(Y^{n-2})\in\mathrm{Cat},
\]
and higher categorical objects to lower-dimensional strata, in a target
\((\infty,n)\)-category chosen before the assignment is made.  Dualizability
conditions ensure that gluing along all strata is represented by composition
and evaluation maps.  In a QFT realization, the target objects must be
constructed from regulated state spaces, boundary operator algebras, and
defect junction spaces.

\begin{openproblem}[Gauge-theoretic extended construction]
\label{op:gauge-theoretic-extended-tqft-construction}
Construct the extended boundary and defect categories of four-dimensional
cohomological gauge theories directly from regulated BV-BFV data, including
instanton and monopole boundary strata, line and surface defects, gluing
maps, orientations, and anomaly lines.
\end{openproblem}
